\documentclass{article}
\usepackage[french]{babel}
\usepackage[utf8]{inputenc}
\usepackage{enumitem}
\setlist[itemize]{label=\textbullet, leftmargin=*}
\usepackage{geometry}
\geometry{a4paper, margin=1in}

\title{Système sécurisé de gestion et de vérification de certificats basé sur la blockchain}
\author{}
\date{February 7, 2026}

\begin{document}

\maketitle

\section{Titre du Projet}
SmartCert
\section{Problématique \& Contexte}
\begin{itemize}
    \item \textbf{Besoin réel :} Les institutions éducatives souhaitent disposer d'un système fiable et sécurisé pour émettre et vérifier des certificats numériques (diplômes, attestations de formation) sans recourir à des démarches administratives lourdes et sujettes aux erreurs ou à la fraude.
    \item \textbf{Public cible :} établissements d'enseignement, étudiants et personnel administratif chargé de l'émission des certificats.
    \item \textbf{Solution numérique :} L'utilisation de la technologie blockchain permet de garantir l'authenticité, la traçabilité et l'immuabilité des certificats, offrant une vérification instantanée par tout tiers sans dépendance aux services administratifs.
    \item \textbf{Exemples concrets :} Réduction des délais de vérification, élimination des risques de falsification, simplification des échanges entre institutions et bénéficiaires.
\end{itemize}

\section{Objectifs du Projet}
\begin{itemize}
    \item \textbf{Buts principaux :} Développer une application web permettant l'émission sécurisée de certificats numériques et leur vérification instantanée via la blockchain, avec un processus complet d'authentification et de stockage.
    \item \textbf{Changements apportés :} Suppression de la dépendance aux processus administratifs manuels, réduction des risques de fraude, accélération des vérifications et amélioration de l'expérience      utilisateur.
    \item \textbf{Différences avec les solutions existantes :} Utilisation de la blockchain pour l'immuabilité des données (contrairement aux systèmes centralisés), avec un mécanisme de vérification transparent et accessible à tout tiers.
    \item \textbf{Valeur ajoutée :} Solution pratique, technique et pédagogique qui apporte une valeur réelle aux institutions éducatives.
\end{itemize}

\section{Description de la Solution \& Valeur Ajoutée}
\begin{itemize}
    \item \textbf{Idée globale :} Création d'une application web intégrant une blockchain (Ethereum Testnet) pour l'émission et la vérification de certificats numériques sécurisés, avec une interface intuitive accessible via navigateur.
    \item \textbf{Fonctionnalités principales :} 
    \begin{itemize}
        \item Émission de certificats avec identifiant unique et cachet numérique
        \item Vérification en temps réel via un lien de vérification
        \item Gestion des certificats via une interface web sécurisée
        \item Enregistrement automatique sur la blockchain via MetaMask
    \end{itemize}
    \item \textbf{Fonctionnalités secondaires :} 
    \begin{itemize}
        \item Téléchargement du certificat au format PDF
        \item Envoi d'email automatique au bénéficiaire
    \end{itemize}
    \item \textbf{Innovation :} Combinaison de technologies web (Flask) et blockchain (Web3.py) pour une solution robuste, sécurisée et évolutive .
    \item \textbf{Technologies utilisées :} 
    \begin{itemize}
        \item Frontend: HTML + CSS + JavaScript...
        \item Backend: Python + Flask
        \item Base de données: SQLite
        \item Blockchain: Ethereum Testnet + Web3.py
        \item Portefeuille: MetaMask

    \end{itemize}
\end{itemize}
\section{Intégration d’un Chatbot via n8n}

Afin d’améliorer l’interaction avec les utilisateurs et d’automatiser certaines tâches administratives, le projet intègre un chatbot intelligent basé sur la technologie n8n.

\begin{itemize}
    \item \textbf{Rôle du chatbot :}
    \begin{itemize}
        \item Assistance automatisée aux utilisateurs (étudiants et administration)
        \item Réponses aux questions fréquentes sur les certificats
        \item Vérification guidée d’un certificat via son identifiant
    \end{itemize}

    \item \textbf{Fonctionnement technique :}
    \begin{itemize}
        \item Le chatbot est implémenté à l’aide de workflows n8n
        \item Communication avec l’application Flask via des API REST

    \end{itemize}

    \item \textbf{Apport au projet :}
    \begin{itemize}
        \item Réduction de la charge du support administratif
        \item Accessibilité 24/7 pour les utilisateurs
        \item Amélioration de l’expérience utilisateur
        \item Automatisation fiable des processus liés aux certificats
    \end{itemize}
\end{itemize}

\end{document}